\documentclass[../../../include/open-logic-section]{subfiles}

\begin{document}

\olfileid{sth}{story}{extensionality}
\olsection{ઘટકો દ્વારા નિર્ધારિત સમાનતા}

સૌપ્રથમ કહેવાની વાત એ છે કે ગણોની ઓળખ તેમના !!{element}s પરથી
નક્કી થાય છે. વધુ ચોક્કસ રીતે:

\begin{axiom}[ઘટકો દ્વારા નિર્ધારિત સમાનતા]
જો ગણો $A$ અને $B$ના !!{element}s એકના એક જ હોય, તો $A$ અને $B$
એક જ ગણ છે.
\[
  \lforall[A][\lforall[B][(\lforall[x][(x \in A \liff x \in B)] \lif
  \eq[A][B])]]
\]
\end{axiom}

\olref[sfr][][]{part}માં આપણે આ વાત સતત માની હતી. છતાં તેને ફરી
કહેવી યોગ્ય છે. ઘટકો દ્વારા નિર્ધારિત સમાનતાની સ્વયંસિદ્ધિ એ
મૂળભૂત વિચાર વ્યક્ત કરે છે કે ગણ તેના !!{element}s દ્વારા નક્કી
થાય છે. (આ રીતે ગણોને \emph{સંકલ્પનાઓ}થી જુદા પાડી શકાય: એકની એક
વસ્તુઓ ઘણી જુદી સંકલ્પનાઓ હેઠળ આવી શકે છે.)

આ સિદ્ધાંત શા માટે સ્વીકારવો? ઘટકો દ્વારા નિર્ધારિત સમાનતાનો કોઈ પણ
ઇનકાર, જેને \emph{ગણસિદ્ધાંત} કહી શકાય તેવું બધું જ છોડી દેવાનો
નિર્ણય છે, એમ કહેવું વાજબી લાગે છે. ગણસિદ્ધાંત એટલે ઘટકો દ્વારા
નિર્ધારિત સંગ્રહોનો જ સિદ્ધાંત.

પરંતુ \olref[sth][][]{part}માં ખરેખરો પડકાર એવા સિદ્ધાંતો નક્કી
કરવાનો છે જે આપણને કહે કે \emph{કયા ગણો અસ્તિત્વ ધરાવે છે}. અને
આ પ્રશ્નનો એકમાત્ર ખરેખર ``સ્વાભાવિક'' લાગતો જવાબ ખોટો છે, એ
સાબિત કરી શકાય છે.

\end{document}
